\documentclass{article}
\usepackage[utf8]{inputenc}
\usepackage[T1]{fontenc}
\usepackage[margin=2cm]{geometry}
\usepackage{hyperref}
\sloppy
\begin{document}
\section*{Historical source}
Encyclopaedia Britannica scan, source row 266. Parse 3.5 extracted the text. Super 3.5 VL corrected it from an enlarged crop of each block. It can still contain recognition errors.
Page headers, footers, and page numbers are removed.

\subsection*{Untitled continuation (E0017)}
\textit{Continues an entry from an earlier page.}

art III.

4. Marriage is forbidden within certain degrees of blood. By the law of Moses (Leviticus xviii.), which, by the act 1567, c. 15. has been adopted by us, seconds in blood, and all remote degrees, may all lawfully marry. By seconds in blood are meant first cousins. Marriage in the direct line is forbidden in infinitum; as it is also in the collateral line in the special case where one of the parties is loco parentis to the other, as grand-uncle, great grand-uncle, \&c. with respect to his grand niece, \&c. The same degrees that are prohibited in consanguinity, are prohibited in affinity; which is the tie arising from marriage, betwixt one of the married pair, and the blood relations of the other. Marriage also, where either of the parties is naturally unfit for generation, or stands already married to a third person, is ipso jure null.

5. To prevent bigamy and incestuous marriages, the church has introduced proclamation of banns; which is the ceremony of publishing the names and designations of those who intend to intermarry, in the churches where the bride and bridegroom reside; after the congregation is assembled for divine service; that all persons who know any objection to the marriage may offer it. When the order of the church is observed, the marriage is called regular; when otherwise, clandestine. Marriage is valid when entered into in either of these ways; but when clandestine, there are certain penalties imposed upon the parties as well as the celebrant and witnesses.

6. By marriage, a society is created between the married pair, which draws after it a mutual communication of their civil interests, in as far as is necessary for maintaining it. As the society lasts only for the joint lives of the socii; therefore rights that have the nature of a perpetuity, which our law flies heritable, are not brought under the partnership or communion of goods; as a land-ellate, or bonds bearing a yearly interest: it is only movable subjects, or the fruits produced by heritable subjects during the marriage, that become common to man and wife.

7. The husband, as the head of the wife, has the sole right of managing the goods in communion, which is called jus mariti. This right is so absolute, that it bears but little resemblance to a right of administering a common object. For the husband can, in virtue thereof, sell, or even gift, at his pleasure, the whole goods falling under communion; and his creditors may affect them for the payment of his proper debts: so that the jus mariti carries all the characters of an affiliation, by the wife to her husband, of her moveable estate. It arises ipso jure from the marriage; and therefore needs no other condition. But a stranger may convey an estate to a wife, so as it shall not be subject to the husband's administration; or the husband himself may, in the marriage-contract, renounce his jus mariti in all or any part of his wife's moveable estate.

8. From this right are excepted paraphernal goods, which, as the word is understood in our law, comprehends the wife's wearing apparel, and the ornaments proper to her person; as necklaces, ear-rings, breast or arm jewels, buckles, \&c. These are neither alienable by the husband, nor affectable by his creditors. Things of promiscuous use to husband and wife, as plate, medals, \&c. may become paraphernal, by the

husband's giving them to the wife, at or before mar-
riage; but they are paraphernal only in regard to
that husband who gave them as such, and are esteemed
common movables, if the wife, whose parapherna-
lia they were, be afterwards married to a second hus-
band; unless he fall in the same manner appropriate
them to her.

9. The right of the husband to the wife's moveable estate, is burdened with the moveable debts contracted by her before marriage: and as his right is universal, so also is his burden; for it reaches to her whole moveable debts, though they should far exceed her moveable estate. Yet the husband is not considered as the true debtor in his wife's debts. In all actions for payment, she is the proper defender: the husband is only cited for his interest, that is, as curator to her, and administrator of the fœcundity-goods. As soon therefore as the marriage is dissolved, and the fœcundity goods thereby (after a division, the husband is no farther concerned in the share belonging to his deceased wife; and consequently is no longer liable to pay her debts, which may be recovered from her representatives, or her separate estate.

10. This obligation upon the husband is, however, perpetuated against him (1.) Where his proper estate, real or personal, has been affected, during the marriage, by complete legal diligence; in which case, the husband must, by the common rules of law, relieve his property from the burden with which it stands charged; but the utmost diligence against his person is not sufficient to perpetuate the obligation; nor even incomplete diligence against his estate. (2.) The husband continues liable, even after the wife's death, in so far as he is lucratus or profited by her estate: Still, however, the law does not consider a husband who has got but a moderate tocher with the wife as lucratus by the marriage; it is the excess only which it considers as lucrum, and that must be estimated by the quality of the parties and their condition of life.—As he was at no time the proper debtor in his wife's moveable debts; therefore, though he should be lucratu, he is, after the diffolation, only liable for them fubfidiariae, i. e. if her own separate estate is not sufficient to pay them off.

11. Where the wife is debtor in that sort of debt, which, if it had been due to her, would have excluded the jus mariti, e.g. in bonds bearing interest, which, as we shall afterwards fee (clxiii. 4.), continues heritable as to the rights of husband and wife, notwithstanding of the enactment of the statute 1661, which renders them moveable in certain other respects, the husband is liable only for the bygone interest, and those that may grow upon the debt during the marriage; because his obligation for her debts must be commensurated to the interest he has in her estate. It is the husband alone who is liable in personal diligence for his wife's debts, while the marriage subsists: the wife, who is the proper debtor, is free from all personal execution upon them while she is uxore viro.

12. The husband by marriage becomes the perpetual curator of the wife. From this right it arises, and is the title "must endure" eug.


\section*{Source record}
\begin{description}
\item[Scan ID] \texttt{britannica:309f765b1ecc83e0a7f7762e44b4a927a7555443:train:266}
\item[Image file] \texttt{page-0266-dfbe51d09127.jpg}
\item[Image SHA-256] \texttt{dfbe51d091279145e37f51fdee7d7b6903c89ea6d64f90fb27913a7d0437e246}
\item[Dataset] NationalLibraryOfScotland/encyclopaedia\_britannica\_illustrated at revision \texttt{309f765b1ecc83e0a7f7762e44b4a927a7555443} (CC0-1.0)
\item[Extraction] model \texttt{nvidia/nemotron-3.5-parse}, raw reply \texttt{extract/raw/page-0266-dfbe51d09127.json}
\item[Raw reply SHA-256] \texttt{a03196b57844b1d3b159caaece105823ce9c0a6582d0f26ee3e0c74c16b56166}
\item[Correction] model \texttt{nvidia/Nemotron-3.5-Super-VL}, block crops, raw replies \texttt{correct/raw/crop\_nothink/page-0266-dfbe51d09127-b*.json}; super\_vl: 11
\item[Running header] none
\end{description}
\end{document}
